La metodología es el producto del trabajo. Esta sección la presenta tal como resulta del análisis comparado: sus requisitos, sus fases y pasos, sus herramientas de diagnóstico, su modelo de estimación y sus reglas de decisión, su catálogo de escenarios y sus procedimientos. Cada fase se vincula con la obligación de la LOPDP que atiende y con el referente internacional del que se toma la buena práctica; los elementos que dependan de la síntesis comparada se ajustarán con los resultados de la Sección~\ref{sec:res_comparado} y con las observaciones del panel Delphi (Sección~\ref{sec:res_delphi}).

\subsection{Descripción de la Propuesta}
\label{sec:descripcion_propuesta}

La propuesta es una \textbf{metodología de diagnóstico y gestión de riesgos en el tratamiento de datos personales académicos} dirigida a cualquier universidad ecuatoriana, compuesta por cuatro productos:

\begin{enumerate}
  \item \textbf{Documento metodológico}: fases P y 0 a 4, escalas, tabla de decisión, criterios de aceptación y procedimiento de respuesta a vulneraciones (esta sección).
  \item \textbf{Herramientas de diagnóstico}: el perfil institucional (Anexo~\ref{anx:perfil}), la lista de verificación de requisitos (Tabla~\ref{tab:lista_verificacion}), la escala de madurez (Anexo~\ref{anx:madurez}), la tabla de finalidades y bases de legitimación (Anexo~\ref{anx:finalidad_base}) y la lista de supuestos de evaluación de impacto obligatoria (Anexo~\ref{anx:eipd}).
  \item \textbf{Plantilla reutilizable de matriz de riesgos}: archivo \path{matriz/Plantilla_Matriz_Riesgos_LOPDP.xlsx}, con quince hojas ---entre ellas, la normativa de referencia, el perfil institucional, el inventario de encargados y transferencias, la revisión previa de EIPD y la tabla de retención---, listas desplegables, fórmulas y formato condicional (Anexo~\ref{anx:plantilla_matriz}).
  \item \textbf{Matriz de escenarios sectoriales}: archivo \path{matriz/Matriz_escenarios_sectoriales.xlsx}, con los quince escenarios de la Sección~\ref{sec:escenarios_sectoriales}, generado desde \path{matriz/escenarios_sectoriales.json}.
  \item \textbf{Guía de aplicación} (Anexo~\ref{anx:guia_aplicacion}): versión operativa de la metodología, con roles, pasos numerados, entradas y productos, que una universidad puede seguir sin leer el resto de la tesis.
\end{enumerate}

La plantilla y la matriz se generan con un programa en JavaScript (\path{scripts/build_matriz.js}) que construye el libro de cálculo, evalúa todas sus fórmulas con un motor de cálculo independiente y reporta los resultados, de modo que las cifras que se trasladen al documento procedan de las fórmulas de la matriz y no de cálculos manuales.

\subsection{Requisitos de la Metodología}
\label{sec:requisitos_guia}

Dado que el artefacto es una metodología y no un sistema informático, los requisitos se expresan como condiciones que la metodología debe satisfacer. Los funcionales describen lo que debe permitir hacer y se derivan de la LOPDP y de la guía de la SPDP; los no funcionales describen las cualidades que debe tener para que una universidad la adopte.

\subsubsection{Requisitos funcionales}

La Tabla~\ref{tab:rf} enumera los requisitos funcionales y su fundamento normativo.

\begin{table}[htbp]
  \centering
  \caption{Requisitos funcionales de la metodología.}
  \label{tab:rf}
  \footnotesize
  \renewcommand{\arraystretch}{1.2}
  \begin{tabularx}{\textwidth}{|c|L|p{3.6cm}|}
    \hline
    \textbf{ID} & \textbf{Requisito} & \textbf{Fundamento} \\
    \hline
    RF01 & Diagnosticar la transparencia hacia el titular y la madurez de la institución antes de iniciar la gestión de riesgos, con el perfil institucional, la lista de verificación de requisitos y la escala de madurez de la SPDP. & LOPDP arts.~12, 37 y 39; Res. SPDP-SPD-2025-0040-R; fase 0 \\
    \hline
    RF02 & Inventariar las actividades de tratamiento (RAT) con finalidades, categorías de datos, titulares, destinatarios, encargados y plazos de conservación. & LOPDP art.~40; guía SPDP (contexto) \\
    \hline
    RF03 & Identificar vulnerabilidades de cuatro tipos: jurídicas, organizacionales, técnicas y relativas a grupos vulnerables de titulares. & Guía SPDP \\
    \hline
    RF04 & Construir escenarios de riesgo con comunidad de amenaza, vector, vulnerabilidades e impacto primario y secundario, desde la perspectiva del titular y del responsable. & Guía SPDP \\
    \hline
    RF05 & Estimar la probabilidad referida a un lapso anual y el impacto, con justificación documentada de cada valor. & Guía SPDP \\
    \hline
    RF06 & Modelar por separado la confidencialidad, la integridad y la disponibilidad, además de la licitud del tratamiento. & Guía SPDP; LOPDP art.~37 \\
    \hline
    RF07 & Ponderar el impacto con los criterios de categorías especiales, grupos vulnerables, número de titulares, naturaleza de la vulneración y volumen de datos. & LOPDP arts.~25 y 40; guía SPDP \\
    \hline
    RF08 & Determinar el nivel de riesgo inherente y residual con una regla explícita para los escenarios de baja probabilidad y muy alto impacto. & Guía SPDP \\
    \hline
    RF09 & Determinar la necesidad de una evaluación de impacto (EIPD) con una lista de supuestos obligatorios aplicada al RAT antes del análisis de riesgos, y además cuando la matriz arroje un nivel alto o muy alto (RA-4). & LOPDP art.~42; Res. SPDP-SPD-2026-0005-R \\
    \hline
    RF10 & Seleccionar opciones de tratamiento (mitigar, transferir, evitar, aceptar) y controles con responsable, plazo e indicador. & LOPDP art.~37; ISO/IEC 27005 \\
    \hline
    RF11 & Definir indicadores de monitoreo y un procedimiento de respuesta que permita notificar vulneraciones en el término legal. & LOPDP arts.~37 y 43 \\
    \hline
    RF12 & Aplicar la protección de datos desde el diseño y por defecto a los tratamientos nuevos y a los cambios en los existentes, con un control en la matriz (revisión previa del tratamiento y minimización por defecto). & LOPDP arts.~39 y 47, num.~9; Res. SPDP-SPD-2025-0040-R \\
    \hline
    RF13 & Producir en cada fase evidencia documentada (actas, registros, justificaciones y aprobaciones) que permita demostrar el cumplimiento ante la SPDP. & LOPDP arts.~10, lit.~k, y 47, num.~2 \\
    \hline
  \end{tabularx}
  \raggedright
  \elaboradopor{Ramón Alejandro Manzaba Castro}
  \fuente{Elaboración propia con base en \cite{lopdp2021,spdp2026guiav2,spdp2026granescala,spdp2025diseno,iso27005_2022}}
\end{table}

\subsubsection{Requisitos no funcionales}

La Tabla~\ref{tab:rnf} presenta los requisitos no funcionales.

\begin{table}[htbp]
  \centering
  \caption{Requisitos no funcionales de la metodología.}
  \label{tab:rnf}
  \footnotesize
  \renewcommand{\arraystretch}{1.2}
  \begin{tabularx}{\textwidth}{|c|p{2.8cm}|L|}
    \hline
    \textbf{ID} & \textbf{Atributo} & \textbf{Descripción} \\
    \hline
    RNF01 & Trazabilidad & Cada actividad y cada campo de la matriz se vinculan con un objetivo específico, una etapa de la guía de la SPDP y un artículo de la LOPDP. \\
    \hline
    RNF02 & Reutilización & La matriz se entrega como plantilla de hoja de cálculo con listas desplegables, fórmulas y formato condicional, aplicable a cualquier proceso de cualquier universidad. \\
    \hline
    RNF03 & Usabilidad & La metodología puede ser aplicada por personal de TI, jurídico o de gestión sin formación especializada en análisis cuantitativo de riesgos. \\
    \hline
    RNF04 & Escalabilidad & Puede aplicarse a un proceso, a una unidad o a toda la institución sin cambiar su estructura. \\
    \hline
    RNF05 & Bajo costo & Se apoya en normas de acceso público y en herramientas ofimáticas de uso común. \\
    \hline
    RNF06 & Confidencialidad & Los expedientes de aplicación usan datos seudonimizados o agregados y registran el acceso a la evidencia. \\
    \hline
    RNF07 & Mantenibilidad & Las escalas, catálogos y criterios se mantienen en hojas separadas para actualizarlos ante cambios normativos, y la hoja \textit{Normativa} registra la fecha de corte, la versión y la fecha de cada norma. \\
    \hline
    RNF08 & Aplicabilidad nacional & La metodología puede aplicarse a cualquier universidad o escuela politécnica del país, pública o privada, de cualquier tamaño y modalidad, ajustando sus obligaciones mediante el perfil institucional. \\
    \hline
  \end{tabularx}
  \raggedright
  \elaboradopor{Ramón Alejandro Manzaba Castro}
  \fuente{Elaboración propia}
\end{table}



Los requisitos fijan, antes del diseño, el criterio con el que se juzgará si la metodología es adecuada; así, la evaluación de la Sección~\ref{sec:validacion_delphi} no depende de la opinión del autor, sino del cumplimiento verificable de cada requisito.

\subsection{Diseño de la Metodología}
\label{sec:diseno_solucion}

\subsubsection{Estructura general de la metodología}

La metodología se compone de seis fases. La \textbf{fase P (preparación normativa)} fija la normativa de referencia, el mandato y el perfil institucional (pasos D0 a D2); la \textbf{fase 0 (diagnóstico institucional)} levanta el inventario de tratamientos y el nivel de madurez (pasos D3 a D11); la \textbf{fase 1 (gestión de riesgos)} establece los criterios, identifica, analiza, evalúa y trata los riesgos (pasos G1, G3 a G5, G7 y G8); la \textbf{fase 2 (evaluación de impacto)} aplica la revisión previa y realiza las evaluaciones obligatorias (pasos G2 y G6); la \textbf{fase 3 (respuesta a vulneraciones)} establece el procedimiento de notificación y comunicación (paso G9), y la \textbf{fase 4 (monitoreo y mejora continua)} mide los indicadores, revisa la matriz y archiva la evidencia (pasos G10 y G11). La preparación se sitúa al inicio porque tanto ISO/IEC 27035-1:2023 como NIST SP 800-61r3 colocan antes de cualquier evaluación las condiciones que la hacen posible: compromiso de la dirección, roles, criterios y procedimientos \cite{iso27035_1_2023,iso27035_2_2023,nist2025sp80061r3}. Las fases se articulan con las etapas obligatorias de la guía de la SPDP, que sigue el proceso de ISO/IEC 27005:2022 \cite{spdp2026guiav2,iso27005_2022}, y se apoyan en dos actividades transversales: la comunicación y consulta con las partes interesadas, y el registro documental de cada decisión, que materializa la responsabilidad proactiva (RF13). La Figura~\ref{fig:fases_metodologia} representa esta estructura.

\begin{figure}[htbp]
  \centering
  \begin{tikzpicture}[
      font=\scriptsize,
      fase/.style={draw, rounded corners=3pt, fill=gray!12, align=center, text width=1.85cm, minimum height=1.55cm, inner sep=3pt},
      trans/.style={draw, align=center, text width=13.9cm, minimum height=0.6cm, inner sep=2pt},
      flecha/.style={-{Stealth[length=2mm]}, thick}
    ]
    \node[fase] (FP) at (0,0)      {\textbf{Fase P}\\Preparación normativa\\(D0--D2)};
    \node[fase] (F0) at (2.4,0)    {\textbf{Fase 0}\\Diagnóstico institucional\\(D3--D11)};
    \node[fase] (F1) at (4.8,0)    {\textbf{Fase 1}\\Gestión de riesgos\\(G1, G3--G5, G7, G8)};
    \node[fase] (F2) at (7.2,0)    {\textbf{Fase 2}\\Evaluación de impacto\\(G2, G6)};
    \node[fase] (F3) at (9.6,0)    {\textbf{Fase 3}\\Respuesta a vulneraciones\\(G9)};
    \node[fase] (F4) at (12.0,0)   {\textbf{Fase 4}\\Monitoreo y mejora\\(G10, G11)};
    \foreach \a/\b in {FP/F0,F0/F1,F1/F2,F2/F3,F3/F4} { \draw[flecha] (\a.east) -- (\b.west); }
    \draw[flecha] (F4.north) -- ++(0,0.55) -| node[pos=0.25, above] {Revisión periódica y ante cambios o incidentes} (F1.north);
    \node[trans] at (6.0,-1.45) {Comunicación y consulta con partes interesadas (DPD, TI, jurídico, unidades académicas)};
    \node[trans] at (6.0,-2.2)  {Registro documental y justificación (\textit{rationale}) de cada estimación y decisión};
  \end{tikzpicture}
  \caption{Estructura de la metodología propuesta.}
  \label{fig:fases_metodologia}
  \raggedright
  \elaboradopor{Ramón Alejandro Manzaba Castro}
  \fuente{Elaboración propia con base en \cite{spdp2026guiav2,iso27005_2022}}
\end{figure}

La Tabla~\ref{tab:fases_detalle} detalla, para cada fase, el objetivo y las actividades, los productos y la evidencia que queda, los roles, el fundamento en la LOPDP y el referente internacional del que se toma la buena práctica. La equivalencia de cada fase con ISO 31000/ISO/IEC 27005, ISO/IEC 27035-1:2023 y NIST SP 800-61r3 se presentó en la Tabla~\ref{tab:equivalencia_marcos}.

{\scriptsize
\renewcommand{\arraystretch}{1.15}
\begin{xltabular}{\textwidth}{|>{\raggedright\arraybackslash}p{1.6cm}|L|>{\raggedright\arraybackslash}p{2.5cm}|>{\raggedright\arraybackslash}p{1.9cm}|>{\raggedright\arraybackslash}p{1.5cm}|>{\raggedright\arraybackslash}p{2.4cm}|}
  \caption{Fases de la metodología: objetivo, productos, roles y referentes.}
  \label{tab:fases_detalle} \\
  \hline
  \textbf{Fase} & \textbf{Objetivo y actividades} & \textbf{Productos y evidencia} & \textbf{Roles} & \textbf{LOPDP} & \textbf{Referente internacional} \\
  \hline
  \endfirsthead
  \hline
  \textbf{Fase} & \textbf{Objetivo y actividades (continuación)} & \textbf{Productos y evidencia} & \textbf{Roles} & \textbf{LOPDP} & \textbf{Referente internacional} \\
  \hline
  \endhead
  \hline
  \multicolumn{6}{r}{\textit{Continúa en la página siguiente}} \\
  \endfoot
  \hline
  \endlastfoot
  P. Preparación normativa & Fijar la normativa de referencia y la fecha de corte; obtener el mandato de la máxima autoridad; definir alcance y roles; responder el perfil institucional (D0 a D2). & Hoja de normativa; acta de inicio; ficha de perfil con obligaciones activadas & Máxima autoridad; delegado; jurídico & Arts.~37, 47 y 48 & Responsabilidad proactiva (RGPD, arts.~5 y 24) \cite{rgpd2016} \\ \hline
  0. Diagnóstico institucional & Mapear procesos y titulares; levantar el RAT con el mapa de datos sensibles; trazar el ciclo de vida; aplicar la lista de verificación; asignar el nivel de madurez; analizar brechas; emitir el informe (D3 a D11). & RAT; diagramas de flujo; lista de verificación calificada; nivel de madurez; informe de diagnóstico & Responsables de proceso; TI; jurídico & Arts.~10, lit.~k, 12, 37 y 39 & Registro de actividades (RGPD, art.~30); escala de madurez de la SPDP \cite{rgpd2016,spdp2025diseno} \\ \hline
  1. Gestión de riesgos & Fijar escalas y criterios; identificar escenarios; estimar probabilidad e impacto con justificación; evaluar con la tabla de decisión; elegir controles; aprobar el riesgo residual (G1, G3 a G5, G7, G8). & Criterios aprobados; matriz inherente y residual; plan de tratamiento; actas de aceptación & Responsables de proceso; TI; panel de estimación & Arts.~37, 40 y 47, num.~5 & Enfoque basado en riesgos (RGPD, art.~24); guía y herramienta Gestiona de la AEPD \cite{aepd2021riesgo,aepd2024gestiona} \\ \hline
  2. Evaluación de impacto & Aplicar la lista de supuestos de evaluación obligatoria antes del análisis; realizar cada EIPD con la plantilla mínima; consultar a la autoridad si el riesgo residual sigue alto (G2, G6). & Lista de EIPD obligatorias; informes de EIPD con la opinión del delegado & Delegado; jurídico; responsables de proceso & Art.~42 & Supuestos y consulta previa (RGPD, arts.~35 y 36); informe de impacto (LGPD, art.~38) \cite{rgpd2016,brasil2018lgpd} \\ \hline
  3. Respuesta a vulneraciones & Detectar, contener, clasificar, notificar a la SPDP y comunicar a los titulares dentro de los términos legales; analizar lecciones (G9). & Registro de incidentes; notificaciones en el SISPDP; comunicaciones a titulares & TI; delegado; jurídico & Arts.~43 y 46 & Notificación de violaciones (RGPD, arts.~33 y 34); incidentes (Res. CD/ANPD 15/2024) \cite{rgpd2016,anpd2024incidentes} \\ \hline
  4. Monitoreo y mejora continua & Medir los indicadores; revisar la matriz al cierre de cada periodo y ante cambios o incidentes; archivar la evidencia (G10, G11). & Tablero de indicadores; bitácora de revisiones; expediente de cumplimiento & Delegado; responsables de proceso & Arts.~10, lit.~k, y 47, num.~3 & Revisión periódica de medidas (RGPD, art.~24) \cite{rgpd2016} \\ \hline
\end{xltabular}
}
{\raggedright\elaboradopor{Ramón Alejandro Manzaba Castro}\fuente{Elaboración propia con base en \cite{lopdp2021,spdp2026guiav2,iso27005_2022} y en los marcos comparados}\par}
% TODO: confirmar los referentes internacionales de cada fase con la síntesis comparada de la Sección 4.1 al cerrar la fase 2 del estudio.

\subsubsection{Pasos de la metodología}
\label{sec:pasos_metodologia}

La metodología se ejecuta en dos secuencias de pasos que aplica la universidad que la adopta. La Tabla~\ref{tab:pasos_diagnostico} presenta los doce pasos del diagnóstico (D0 a D11), que conforman la fase~0 y terminan en una línea base, y la Tabla~\ref{tab:pasos_gestion}, los once pasos de la gestión de riesgos (G1 a G11), que empiezan con esa línea base y se repiten en cada revisión. La última columna de cada tabla indica qué hace este trabajo con cada paso, porque, al no haberse aplicado la metodología en una institución real, algunos pasos se diseñan como plantilla o como lista de verificación y otros se describen para su aplicación por cada universidad.

\begin{table}[htbp]
  \centering
  \caption{Pasos de la metodología de diagnóstico (fase 0).}
  \label{tab:pasos_diagnostico}
  \scriptsize
  \renewcommand{\arraystretch}{1.12}
  \begin{tabularx}{\textwidth}{|c|>{\raggedright\arraybackslash}p{3.0cm}|>{\raggedright\arraybackslash}p{3.4cm}|L|}
    \hline
    \textbf{Paso} & \textbf{Objetivo} & \textbf{Producto} & \textbf{En este trabajo} \\
    \hline
    D0 & Fijar la normativa de referencia & Hoja \textit{Normativa} con la fecha de corte & Corpus de la tesis como punto de partida (Tabla~\ref{tab:corpus_normativo}; Anexo~\ref{anx:corpus}) \\ \hline
    D1 & Mandato y alcance & Acta de inicio y plan de trabajo & No se ejecuta; se describe en la guía de aplicación \\ \hline
    D2 & Perfil institucional & Ficha de perfil y lista de obligaciones activadas & Se diseña como lista de verificación (Anexo~\ref{anx:perfil}) y hoja de la plantilla \\ \hline
    D3 & Mapa de procesos y titulares & Mapa de procesos y titulares & Se describe en la guía de aplicación \\ \hline
    D4 & Registro de actividades de tratamiento (RAT) & RAT preliminar, con el mapa de datos sensibles y la fila del aula virtual & Fila tipo del aula virtual y plantilla \\ \hline
    D5 & Ciclo de vida de los datos & Diagramas de flujo de datos & Diagramas de ejemplo en la guía de aplicación \\ \hline
    D6 & Matriz de requisitos e indicadores & Lista de verificación con criterios S/P/N & Se diseña a partir del análisis comparado (Tabla~\ref{tab:lista_verificacion}) \\ \hline
    D7 & Instrumentos & Instrumentos validados & Se diseñan y se someten al juicio de expertos \\ \hline
    D8 & Levantamiento de evidencia & Expediente de evidencias & Se describe en la guía de aplicación \\ \hline
    D9 & Calificación y nivel de madurez & Índice de cumplimiento y nivel de madurez & Escala y regla de asignación definidas (Anexo~\ref{anx:madurez}) \\ \hline
    D10 & Análisis de brechas & Matriz de brechas priorizada & Se describe; plantilla de la matriz \\ \hline
    D11 & Informe de diagnóstico & Informe y línea base, entrada de la fase~1 & Estructura del informe en la guía de aplicación \\ \hline
  \end{tabularx}
  \raggedright
  \elaboradopor{Ramón Alejandro Manzaba Castro}
  \fuente{Elaboración propia}
\end{table}

\begin{table}[htbp]
  \centering
  \caption{Pasos de la metodología de gestión de riesgos (fases 1 a 4).}
  \label{tab:pasos_gestion}
  \scriptsize
  \renewcommand{\arraystretch}{1.12}
  \begin{tabularx}{\textwidth}{|c|>{\raggedright\arraybackslash}p{3.0cm}|>{\raggedright\arraybackslash}p{3.4cm}|L|}
    \hline
    \textbf{Paso} & \textbf{Objetivo} & \textbf{Producto} & \textbf{En este trabajo} \\
    \hline
    G1 & Contexto y criterios & Documento de criterios aprobado & Se diseña y se calibra con el panel Delphi \\ \hline
    G2 & Revisión previa de EIPD & Lista de tratamientos con EIPD obligatoria & Lista de supuestos (Anexo~\ref{anx:eipd}) y hoja \textit{EIPD previa} \\ \hline
    G3 & Identificación de escenarios & Catálogo de escenarios activados & Catálogo de quince escenarios con criterios de activación ajustables \\ \hline
    G4 & Análisis (probabilidad e impacto) & Matriz de riesgo inherente & Escalas, agravantes y estimación Delphi calibrada como técnica \\ \hline
    G5 & Evaluación & Riesgos priorizados & Tabla de decisión y plantilla (Tabla~\ref{tab:matriz_nivel}) \\ \hline
    G6 & Evaluación de impacto (EIPD) & Informe de EIPD & Plantilla mínima de EIPD (Anexo~\ref{anx:guia_aplicacion}) \\ \hline
    G7 & Tratamiento del riesgo & Plan de tratamiento & Controles de referencia (Tabla~\ref{tab:plan_tratamiento}) \\ \hline
    G8 & Riesgo residual y aceptación & Matriz residual y actas & Reglas de aceptación RA-1 a RA-6 \\ \hline
    G9 & Respuesta a vulneraciones & Procedimiento y registro de incidentes & Procedimiento documentado y validado por jueces \\ \hline
    G10 & Monitoreo y mejora & Tablero de indicadores y bitácora & Indicadores definidos; su aplicación es trabajo futuro \\ \hline
    G11 & Comunicación y archivo de evidencia & Expediente de cumplimiento & Trazabilidad de la matriz (RNF01) y evidencia por fase (RF13) \\ \hline
  \end{tabularx}
  \raggedright
  \elaboradopor{Ramón Alejandro Manzaba Castro}
  \fuente{Elaboración propia}
\end{table}

\subsubsection{Fase 0: perfil institucional y nivel de madurez}
\label{sec:fase0_madurez}

La fase~0 empieza por el \textbf{perfil institucional} (paso D2), porque las obligaciones de una universidad no son las mismas para todas: una pública debe designar un delegado \cite[art.~48, num.~1]{lopdp2021}, está sujeta a la transparencia activa \cite[arts.~8, lit.~g, y 19]{lotaip2023} y responde con un régimen sancionatorio propio ---multas de 1 a 10 o de 10 a 20 salarios básicos unificados a los servidores o funcionarios responsables, según la infracción sea leve o grave---, mientras que una particular responde con multas calculadas sobre su volumen de negocio \cite[arts.~71 y 72]{lopdp2021}. La lista de verificación del perfil (Anexo~\ref{anx:perfil}) pregunta el sostenimiento, el tamaño frente al umbral de gran escala, el tratamiento de datos de salud, discapacidad o biometría, la admisión de menores de edad, la modalidad en línea, el uso de servicios en la nube fuera del país y la existencia de sedes, y cada respuesta afirmativa activa obligaciones y reglas de la metodología. Para las universidades públicas, la regla de minimización se aplica también a lo que publican por transparencia: la LOTAIP considera información confidencial los datos personales cuya difusión requiere el consentimiento de su titular y obliga a los sujetos obligados a proteger los datos personales que custodian \cite[arts.~4, num.~5, y 9]{lotaip2023}, de modo que una lista de admitidos o de becarios se publica sin la cédula completa \cite[art.~10, lit.~e]{lopdp2021}.

La fase~0 entrega, además, un \textbf{nivel de madurez} basado en la escala de la Guía de Protección de Datos Personales desde el Diseño y por Defecto de la SPDP \cite{spdp2025diseno}, que la metodología adopta tal como la describen sus análisis publicados \cite{avl2025diseno,tzvs2025diseno} y aplica a nueve bloques de control (Tabla~\ref{tab:madurez}): gobernanza y responsabilidad proactiva; RAT e inventario; transparencia y derechos; bases de legitimación y consentimiento; seguridad; encargados y transferencias; conservación; riesgos, EIPD y protección desde el diseño, y respuesta a vulneraciones. El porcentaje de cada bloque es la proporción de sus controles aplicados con evidencia, y su nivel se asigna con los cortes de la guía; los controles de cada bloque se detallan en el Anexo~\ref{anx:madurez}.

\begin{table}[htbp]
  \centering
  \caption{Escala de madurez de la fase 0.}
  \label{tab:madurez}
  \scriptsize
  \renewcommand{\arraystretch}{1.15}
  \begin{tabularx}{\textwidth}{|>{\raggedright\arraybackslash}p{2.4cm}|>{\raggedright\arraybackslash}p{3.8cm}|L|}
    \hline
    \textbf{Nivel} & \textbf{Descripción} & \textbf{Evidencias que lo sustentan (ejemplos)} \\
    \hline
    0. Caótico & El principio o control no se conoce ni se aplica. & No hay política, responsable ni registro; la ficha documental da N en los indicadores del bloque. \\ \hline
    1. Implícito & Se conoce, pero se aplica en menos del 25\,\% de los casos, sin documentación. & Prácticas aisladas de algunas unidades; sin aprobación formal ni registro. \\ \hline
    2. Temprano explícito & Aplicación parcial y documentada, entre el 25\,\% y el 75\,\% de los casos. & Política aprobada pero incompleta; RAT parcial; contratos con algunos encargados; indicadores con calificación P. \\ \hline
    3. Maduro explícito & Aplicación consolidada, en más del 75\,\% de los casos, con evidencia y revisión periódica. & Política completa y publicada; RAT y contratos completos; registros de revisión; indicadores con calificación S. \\ \hline
  \end{tabularx}
  \raggedright
  \elaboradopor{Ramón Alejandro Manzaba Castro}
  \fuente{Elaboración propia a partir de la escala de \cite{spdp2025diseno}, según \cite{avl2025diseno,tzvs2025diseno}}
\end{table}

\textbf{Regla de asignación del nivel institucional.} El nivel institucional es el \textbf{menor nivel alcanzado en los cuatro bloques críticos} ---gobernanza y responsabilidad proactiva; transparencia y derechos; seguridad, y riesgos, EIPD y protección desde el diseño---, y se informa siempre junto con el perfil completo de los nueve bloques. Se elige el mínimo y no el promedio por dos razones: un promedio permitiría que un bloque maduro compensara otro caótico, cuando la ley exige cumplir cada obligación y no un promedio de ellas, y los bloques críticos son los que se vinculan con infracciones graves, como no implementar medidas de seguridad o no utilizar metodologías de análisis y gestión de riesgos \cite[art.~68]{lopdp2021}. Con ello, la fase~0 sí mide la madurez, por lo que conserva su nombre.


\subsubsection{Lista de verificación de requisitos}
\label{sec:lista_verificacion}

El paso D6 se apoya en una lista de verificación de requisitos derivados de la LOPDP, con la que la universidad califica su propia información y sus documentos en tres niveles: \textbf{S} (el elemento está presente y es completo), \textbf{P} (está presente, pero incompleto o ambiguo) y \textbf{N} (no existe). La lista tiene quince requisitos principales y ocho complementarios (Tabla~\ref{tab:lista_verificacion}); el requisito ID04 se califica como P cuando la institución declara el consentimiento como base de todos sus tratamientos (Anexo~\ref{anx:finalidad_base}), y la revisión del aula virtual (ID17 a ID19) cubre su pantalla de acceso y su contrato con el proveedor. La lista se completará con los elementos que el análisis comparado identifique como mínimos o como buenas prácticas adoptables (Sección~\ref{sec:res_comparado}). % TODO: añadir los requisitos que resulten de la síntesis comparada.

\begin{table}[htbp]
  \centering
  \caption{Lista de verificación de requisitos de la metodología (paso D6).}
  \label{tab:lista_verificacion}
  \scriptsize
  \renewcommand{\arraystretch}{1.15}
  \begin{tabularx}{\textwidth}{|c|L|>{\raggedright\arraybackslash}p{3.6cm}|}
    \hline
    \textbf{ID} & \textbf{Requisito que se verifica en la información y los documentos de la institución} & \textbf{Base legal} \\
    \hline
    ID01 & Existe una política o aviso de privacidad publicado y accesible desde el sitio principal & LOPDP, arts.~10, lit.~c, 12 y 47, num.~4 \\
    \hline
    ID02 & Identifica al responsable del tratamiento con dirección, teléfono y correo electrónico & LOPDP, art.~12, num.~8 \\
    \hline
    ID03 & Declara los fines del tratamiento de los datos de los estudiantes & LOPDP, arts.~10, lit.~d, y 12, num.~1 \\
    \hline
    ID04 & Declara la base de legitimación correcta para cada finalidad (P si usa el consentimiento como base de todo) & LOPDP, arts.~7, 9 y 12, num.~2 \\
    \hline
    ID05 & Declara el tiempo de conservación de los datos & LOPDP, arts.~10, lit.~i, y 12, num.~4 \\
    \hline
    ID06a & Informa las comunicaciones a terceros nacionales (por ejemplo, SENESCYT, IESS, bancos o aseguradoras), con su finalidad y su base & LOPDP, arts.~12, num.~10, y 33 \\
    \hline
    ID06b & Informa las transferencias internacionales, con el país de destino, la garantía y la base & LOPDP, arts.~12, num.~10, y 55 a 59 \\
    \hline
    ID07 & Enumera los derechos del titular (acceso, rectificación y actualización, eliminación, oposición, portabilidad, suspensión del tratamiento y no ser objeto de decisiones automatizadas) y la forma de ejercerlos & LOPDP, arts.~12, nums.~14 y 15, 13 a 17, 19 y 20 \\
    \hline
    ID08 & Ofrece un canal específico y operativo para ejercer derechos (correo, formulario o ventanilla) & LOPDP, arts.~8 y 13 \\
    \hline
    ID09 & Indica dónde y cómo reclamar ante el responsable y ante la Autoridad de Protección de Datos Personales & LOPDP, art.~12, num.~16 \\
    \hline
    ID10 & Informa la posibilidad de revocar el consentimiento & LOPDP, arts.~8 y 12, num.~13 \\
    \hline
    ID11 & Publica la designación del delegado de protección de datos con sus datos de contacto & LOPDP, arts.~12, num.~9, y 48 \\
    \hline
    ID12 & Hace referencia expresa a la LOPDP y a su Reglamento como marco del tratamiento & LOPDP, arts.~10, lit.~a, y 47, num.~1 \\
    \hline
    ID13 & Declara las categorías especiales y los datos sensibles que trata (salud, discapacidad, biométricos, etnia, condición migratoria, identidad de género o nombre social) y su base de legitimación, y por separado los datos socioeconómicos & LOPDP, arts.~4, 25 y 26 \\
    \hline
    ID14 & Los formularios en línea de admisión, inscripción o contacto incluyen una cláusula informativa previa y un consentimiento separado para las finalidades no necesarias para el servicio educativo & LOPDP, arts.~8 y 12 \\
    \hline
    ID15 & Informa sobre el uso de \textit{cookies} u otras tecnologías de seguimiento y permite aceptarlas o rechazarlas & LOPDP, arts.~8 y 12, num.~3 (complementario) \\
    \hline
    ID16 & Indica la fecha de vigencia o de última actualización de la política & LOPDP, art.~10, lit.~c (complementario) \\
    \hline
    ID17 & Identifica la plataforma educativa (aula virtual) y quién la aloja (dominio propio o proveedor) & LOPDP, arts.~12, num.~10, y 34 (complementario) \\
    \hline
    ID18 & Muestra un aviso de privacidad en la pantalla de acceso del aula virtual & LOPDP, art.~12 (complementario) \\
    \hline
    ID19 & Declara el uso de supervisión de exámenes en línea y si emplea biometría & LOPDP, arts.~12, num.~17, 20 y 42 (complementario) \\
    \hline
    ID20 & Describe los tipos de tratamiento, la base de datos en que constan los datos, su origen cuando no se obtuvieron del titular y las finalidades ulteriores & LOPDP, art.~12, nums.~3, 5, 6 y 7 (complementario) \\
    \hline
    ID21 & Informa las consecuencias de entregar o negarse a entregar los datos y el efecto de suministrar datos erróneos & LOPDP, art.~12, nums.~11 y 12 (complementario) \\
    \hline
    ID22 & Informa la existencia de valoraciones o decisiones automatizadas, incluida la elaboración de perfiles & LOPDP, arts.~12, num.~17, y 20 (complementario) \\
    \hline
  \end{tabularx}
  \raggedright
  \elaboradopor{Ramón Alejandro Manzaba Castro}
  \fuente{Elaboración propia a partir de \cite{lopdp2021}}
\end{table}

La Tabla~\ref{tab:cobertura_art12} muestra que cada uno de los diecisiete numerales del artículo~12 de la LOPDP tiene al menos un requisito de la lista que lo verifica.

\begin{table}[htbp]
  \centering
  \caption{Cobertura de los diecisiete numerales del artículo~12 de la LOPDP.}
  \label{tab:cobertura_art12}
  \scriptsize
  \renewcommand{\arraystretch}{1.12}
  \begin{tabularx}{\textwidth}{|c|L|>{\raggedright\arraybackslash}p{4.4cm}|}
    \hline
    \textbf{Num.} & \textbf{Contenido (texto oficial, parafraseado)} & \textbf{Indicador que lo verifica} \\
    \hline
    1 & Los fines del tratamiento & ID03 \\ \hline
    2 & La base legal del tratamiento & ID04 \\ \hline
    3 & Los tipos de tratamiento & ID20; ID15 (seguimiento en línea) \\ \hline
    4 & El tiempo de conservación & ID05 \\ \hline
    5 & La existencia de una base de datos en la que constan sus datos & ID20 \\ \hline
    6 & El origen de los datos cuando no se obtuvieron del titular & ID20 \\ \hline
    7 & Otras finalidades y tratamientos ulteriores & ID20 \\ \hline
    8 & Identidad y datos de contacto del responsable (domicilio, teléfono y correo) & ID02 \\ \hline
    9 & Cuando sea del caso, identidad y contacto del delegado & ID11 \\ \hline
    10 & Transferencias o comunicaciones nacionales o internacionales, destinatarios, finalidades y garantías & ID06a, ID06b e ID17 \\ \hline
    11 & Consecuencias de entregar los datos o de negarse a ello & ID21 \\ \hline
    12 & Efecto de suministrar datos erróneos o inexactos & ID21 \\ \hline
    13 & La posibilidad de revocar el consentimiento & ID10 \\ \hline
    14 & Existencia y forma de ejercer los derechos de acceso, eliminación, rectificación y actualización, oposición, anulación, limitación del tratamiento y a no ser objeto de decisiones automatizadas & ID07 e ID08 \\ \hline
    15 & Mecanismos para ejercer la portabilidad & ID07 \\ \hline
    16 & Dónde y cómo reclamar ante el responsable y la Autoridad & ID09 \\ \hline
    17 & Existencia de valoraciones y decisiones automatizadas, incluida la elaboración de perfiles & ID22 e ID19 \\ \hline
  \end{tabularx}
  \par\vspace{2pt}
  \raggedright
  \textit{Nota:} el momento de la información ---previa, en la recogida, cuando los datos se obtienen del titular--- se verifica con ID14 (formularios) e ID18 (pantalla de acceso del aula virtual).
  \par\vspace{2pt}
  \elaboradopor{Ramón Alejandro Manzaba Castro}
  \fuente{Elaboración propia a partir de \cite[art.~12]{lopdp2021}}
\end{table}
\subsubsection{Modelo de estimación del riesgo}

El modelo de estimación respeta las reglas de la guía de la SPDP: cada escenario se analiza por separado para la confidencialidad, la integridad, la disponibilidad y la licitud; la probabilidad se refiere a un lapso anual; cada valor se acompaña de su justificación, y las escalas se tratan como etiquetas ordinales, no como magnitudes que puedan multiplicarse \cite{spdp2026guiav2}. Las escalas que se proponen a continuación constituyen la versión inicial que el panel Delphi revisará y calibrará.

La Tabla~\ref{tab:escala_probabilidad} define la escala de probabilidad. Cada nivel se asocia con un intervalo de frecuencia anual esperada, lo que obliga a que la estimación tenga un referente verificable y permite contrastarla en el futuro con el registro de incidentes.

\begin{table}[htbp]
  \centering
  \caption{Escala de probabilidad (lapso anual).}
  \label{tab:escala_probabilidad}
  \small
  \renewcommand{\arraystretch}{1.2}
  \begin{tabularx}{\textwidth}{|p{2.2cm}|p{3.6cm}|L|}
    \hline
    \textbf{Nivel} & \textbf{Frecuencia anual esperada} & \textbf{Descripción orientativa} \\
    \hline
    Muy baja & Menos de una vez cada 10 años & No hay antecedentes en la institución ni en el sector; la amenaza requiere capacidades poco comunes. \\
    \hline
    Baja & Una vez cada 2 a 10 años & Hay antecedentes en el sector, pero los controles existentes dificultan su materialización. \\
    \hline
    Media & Aproximadamente una vez cada 1 a 2 años & Hay antecedentes en IES similares y controles parciales. \\
    \hline
    Alta & Entre 1 y 5 veces por año & El evento ha ocurrido en la institución o la vulnerabilidad es conocida y explotable. \\
    \hline
    Muy alta & Más de 5 veces por año & El evento es recurrente o la vulnerabilidad se explota con facilidad. \\
    \hline
  \end{tabularx}
  \raggedright
  \elaboradopor{Ramón Alejandro Manzaba Castro}
  \fuente{Elaboración propia con base en \cite{spdp2026guiav2}}
\end{table}

El impacto se valora en primer lugar desde la perspectiva del titular, conforme al enfoque de derechos de la LOPDP, y en segundo lugar desde la del responsable (sanciones, costos y reputación). La Tabla~\ref{tab:escala_impacto} define la escala de impacto sobre el titular, adaptada de los niveles de severidad propuestos por ENISA \cite{enisa2013severity}.

\begin{table}[htbp]
  \centering
  \caption{Escala de impacto sobre el titular.}
  \label{tab:escala_impacto}
  \small
  \renewcommand{\arraystretch}{1.2}
  \begin{tabularx}{\textwidth}{|p{2.2cm}|L|}
    \hline
    \textbf{Nivel} & \textbf{Consecuencias para el titular} \\
    \hline
    Muy bajo & Molestias menores que el titular supera sin dificultad (por ejemplo, reingresar información). \\
    \hline
    Bajo & Inconvenientes que el titular supera con algún esfuerzo (costos menores, tiempo, molestias). \\
    \hline
    Medio & Consecuencias significativas que el titular puede superar con dificultad (pérdida de oportunidades académicas, cargos indebidos reversibles, temor o estrés). \\
    \hline
    Alto & Consecuencias graves, difíciles de superar (fraude o suplantación de identidad, discriminación, exclusión de beneficios, daño reputacional, exposición de datos socioeconómicos de becas). \\
    \hline
    Muy alto & Consecuencias irreversibles o que afectan derechos fundamentales, la salud o la integridad del titular (exposición de datos de salud o discapacidad, afectación a menores de edad). \\
    \hline
  \end{tabularx}
  \raggedright
  \elaboradopor{Ramón Alejandro Manzaba Castro}
  \fuente{Elaboración propia adaptada de \cite{enisa2013severity}}
\end{table}

El nivel de impacto inicial se ajusta con los criterios agravantes que exige la guía de la SPDP \cite{spdp2026guiav2}: (a) si el escenario involucra categorías especiales de datos o datos de grupos vulnerables, en particular de adolescentes \cite[art.~25]{lopdp2021}, el impacto sube un nivel; (b) si el tratamiento afectado alcanza el umbral de gran escala del Modelo Técnico de Gran Escala de la SPDP o encaja en uno de sus supuestos \cite{spdp2026granescala} ---lo que se verifica en el propio tratamiento con los criterios y umbrales de la norma (Sección~\ref{sec:etapa_ecuador})---, o supera el umbral de titulares que la institución haya definido como masivo, el impacto sube un nivel; (c) si la vulneración es intencional o los datos pueden combinarse para identificar a los titulares con facilidad, el impacto sube un nivel. El ajuste acumulado no puede superar el nivel muy alto y cada ajuste debe justificarse. El impacto se valora siempre desde el titular y según la dimensión afectada: una interrupción de disponibilidad, como la que causa un \textit{ransomware} que solo cifra (R-08), tiene normalmente un impacto medio para el estudiante, que pierde temporalmente el acceso a servicios académicos; la extracción y publicación de sus datos (R-09 y R-15) afecta la confidencialidad y tiene un impacto alto o muy alto cuando incluye datos sensibles, porque se aplica el agravante (a).

\textbf{Fuentes de información de cada estimación.} La guía de la SPDP exige que todo valor tenga una justificación \cite{spdp2026guiav2}; para que esa justificación no dependa del criterio del analista, la Tabla~\ref{tab:fuentes_estimacion} fija de antemano qué fuente respalda cada variable y en qué orden de preferencia se usan.

\begin{table}[htbp]
  \centering
  \caption{Fuentes de información y reglas para cada variable de la estimación.}
  \label{tab:fuentes_estimacion}
  \scriptsize
  \renewcommand{\arraystretch}{1.2}
  \begin{tabularx}{\textwidth}{|>{\raggedright\arraybackslash}p{2.4cm}|L|L|}
    \hline
    \textbf{Variable} & \textbf{Fuentes, en orden de preferencia} & \textbf{Regla de asignación} \\
    \hline
    Probabilidad anual ($P$) & (1) Registro de incidentes de la institución; (2) estado de los controles verificado en la fase 0; (3) brechas del diagnóstico de la fase~0 y antecedentes sectoriales documentados \cite{ulven2021systematic,li2023breach}; (4) estimación Delphi calibrada & Se asigna el nivel cuya frecuencia anual (Tabla~\ref{tab:escala_probabilidad}) contiene la frecuencia estimada; si las fuentes discrepan, prevalece la de mayor preferencia y se registra la discrepancia. Las cifras de encuestas de proveedores comerciales, como \cite{sophos2024education}, se usan solo como contexto y no para estimar frecuencias \\
    \hline
    Impacto sobre el titular ($I_0$) & (1) Categorías de datos y titulares del RAT, incluido el mapa de datos sensibles; (2) criterios de severidad de ENISA \cite{enisa2013severity}; (3) estimación Delphi calibrada & Se asigna el nivel cuyo descriptor (Tabla~\ref{tab:escala_impacto}) corresponde a la consecuencia más grave razonablemente previsible para el titular \\
    \hline
    Agravantes ($a_1, a_2, a_3$) & RAT; volumen de titulares y proporción de menores de edad; Modelo Técnico de Gran Escala \cite{spdp2026granescala} & Cada agravante vale 0 o 1 según se cumpla o no el criterio (a), (b) o (c) \\
    \hline
    Impacto sobre el responsable & Régimen sancionatorio según el perfil institucional: público (multas a servidores) o privado (porcentaje del volumen de negocio) \cite[arts.~71 y 72]{lopdp2021}; costos de respuesta & Se registra como información para la decisión; no modifica el nivel, que se determina por el impacto sobre el titular \\
    \hline
  \end{tabularx}
  \raggedright
  \elaboradopor{Ramón Alejandro Manzaba Castro}
  \fuente{Elaboración propia con base en \cite{spdp2026guiav2,enisa2013severity,spdp2026granescala,lopdp2021}}
\end{table}

\textbf{Estimación por expertos calibrados.} Cuando la institución no dispone de datos históricos suficientes, la estimación se hace con un panel Delphi de expertos propios o externos (Anexo~\ref{anx:guia_aplicacion}, Sección~\ref{sec:guia_delphi}). Para cumplir la exigencia de la guía de la SPDP de validar las estimaciones \cite{spdp2026guiav2}, los expertos responden una batería de preguntas de calibración con respuesta conocida, expresando su grado de confianza como una probabilidad. La calidad de esas respuestas se mide con el puntaje de Brier \cite{brier1950verification}, definido en la Ecuación~\ref{eq:brier}:

\begin{equation}
\label{eq:brier}
BS = \frac{1}{N} \sum_{t=1}^{N} \left( f_t - o_t \right)^2
\end{equation}
\addeqtolist{Puntaje de Brier para la calibración de estimaciones}

\noindent donde $N$ es el número de preguntas, $f_t$ es la probabilidad asignada por el experto a la pregunta $t$ y $o_t$ toma el valor 1 si el evento ocurrió (o la afirmación es verdadera) y 0 en caso contrario. El puntaje varía entre 0, que corresponde a una calibración perfecta, y 1; un experto que respondiera siempre con una probabilidad de 0,5 obtendría 0,25, valor que se adopta como umbral de referencia. Las estimaciones de los expertos calibrados se emplean para ponderar las respuestas del panel conforme a la Ecuación~\ref{eq:ponderacion}, que es la que implementa la hoja \textit{Delphi} de la plantilla:

\begin{equation}
\label{eq:ponderacion}
\hat{x} = \frac{\sum_{e=1}^{m} w_e\, x_e}{\sum_{e=1}^{m} w_e}, \qquad
w_e = \begin{cases} 1 - BS_e & \text{si } BS_e \le 0{,}25 \\ 0 & \text{si } BS_e > 0{,}25 \end{cases}
\end{equation}
\addeqtolist{Estimación del panel ponderada por calibración}

\noindent donde $x_e$ es la estimación ordinal (1 a 5) del experto $e$ y $BS_e$ su puntaje de Brier. La mediana sigue siendo el valor de referencia del consenso; la media ponderada se informa como contraste, y si ambas difieren en más de un nivel, el escenario vuelve a ronda.

Formalmente, el impacto ajustado se obtiene con la Ecuación~\ref{eq:impacto_ajustado}, en la que los niveles se codifican de 1 (muy bajo) a 5 (muy alto) solo para aplicar el tope, no para operar con ellos como magnitudes:

\begin{equation}
\label{eq:impacto_ajustado}
I_{aj} = \min\left(5,\; I_0 + a_1 + a_2 + a_3\right), \qquad a_k \in \{0, 1\}
\end{equation}
\addeqtolist{Impacto ajustado por agravantes}

\subsubsection{Determinación del nivel de riesgo y criterios de aceptación}

Para evitar que el nivel de riesgo se obtenga multiplicando etiquetas, práctica que la guía de la SPDP desaconseja \cite{spdp2026guiav2}, el nivel no se calcula con una fórmula aritmética, sino con una función de decisión $N = T(P, I_{aj})$, donde $T$ es la Tabla~\ref{tab:matriz_nivel}. La tabla otorga mayor peso al impacto que a la probabilidad y contiene una regla explícita para los escenarios de baja probabilidad y alto impacto: todo escenario con impacto muy alto sobre el titular se clasifica, como mínimo, en nivel alto, con independencia de su probabilidad. El mapa de calor resultante se utiliza solo para representar los resultados, no para calibrarlos.

\begin{table}[htbp]
  \centering
  \caption{Tabla de decisión del nivel de riesgo.}
  \label{tab:matriz_nivel}
  \small
  \renewcommand{\arraystretch}{1.35}
  \begin{tabular}{|l|c|c|c|c|c|}
    \hline
    \textbf{Probabilidad / Impacto} & \textbf{Muy bajo} & \textbf{Bajo} & \textbf{Medio} & \textbf{Alto} & \textbf{Muy alto} \\
    \hline
    \textbf{Muy alta} & Medio & Alto & Alto & Muy alto & Muy alto \\
    \hline
    \textbf{Alta}     & Bajo & Medio & Alto & Muy alto & Muy alto \\
    \hline
    \textbf{Media}    & Bajo & Medio & Medio & Alto & Muy alto \\
    \hline
    \textbf{Baja}     & Muy bajo & Bajo & Medio & Alto & Alto \\
    \hline
    \textbf{Muy baja} & Muy bajo & Muy bajo & Bajo & Medio & Alto \\
    \hline
  \end{tabular}
  \par\vspace{2pt}
  \raggedright
  \elaboradopor{Ramón Alejandro Manzaba Castro}
  \fuente{Elaboración propia con base en \cite{spdp2026guiav2}}
\end{table}

Los criterios de aceptación del riesgo se expresan como reglas numeradas, de modo que cualquier revisor pueda verificar su aplicación:

\begin{itemize}
  \item \textbf{RA-1.} Nivel muy bajo o bajo: aceptable, con registro de la decisión y revisión al cierre de cada periodo académico.
  \item \textbf{RA-2.} Nivel medio: se trata dentro del plan anual, con un plazo máximo de doce meses, salvo que la dirección acepte formalmente el riesgo con su justificación.
  \item \textbf{RA-3.} Nivel alto o muy alto: no aceptable; requiere tratamiento prioritario, con plazo máximo de noventa días para iniciar los controles.
  \item \textbf{RA-4.} Nivel alto o muy alto en escenarios con categorías especiales o en tratamientos a gran escala: activa la evaluación de impacto \cite[art.~42]{lopdp2021}. Es un activador \textbf{adicional}: la EIPD obligatoria se determina antes, con la lista de supuestos de la revisión previa (paso G2).
  \item \textbf{RA-5.} Impacto ajustado muy alto ($I_{aj} = 5$): ningún riesgo residual puede aceptarse sin las observaciones previas del delegado de protección de datos y del área jurídica y sin la aprobación de la máxima autoridad.
  \item \textbf{RA-6.} Si, tras la EIPD y los controles propuestos, el riesgo residual de un tratamiento sigue siendo alto o muy alto, el tratamiento no se inicia sin consultar previamente a la SPDP, con el informe de la EIPD y las observaciones del delegado.
\end{itemize}

\textbf{Revisión previa de EIPD obligatoria (paso G2).} Activar la EIPD solo cuando la matriz da un nivel alto y hay categorías especiales o gran escala sería circular: si la estimación fuera baja por error, la EIPD nunca se haría, aunque la ley exige hacerla antes de iniciar el tratamiento \cite[art.~42]{lopdp2021}. Por eso la metodología aplica al RAT, \textbf{antes} del análisis de riesgos, una lista de supuestos en la que un solo supuesto basta para que la EIPD sea obligatoria, con independencia del nivel que después dé la matriz: (1) biometría, como la huella o el reconocimiento facial; (2) supervisión de exámenes en línea con cámara, micrófono o monitoreo; (3) datos de salud, psicológicos o de discapacidad del área de bienestar; (4) perfilamiento o analítica predictiva, como el riesgo de deserción o el rendimiento; (5) tratamiento sistemático de datos de aspirantes menores de edad; (6) tratamientos que alcancen el umbral de gran escala \cite{spdp2026granescala}; (7) tratamientos nuevos o cambios sustanciales en uno existente, que es donde se aplica la protección desde el diseño (RF12); (8) cualquier supuesto que la SPDP haya calificado como de alto riesgo, y (9) la observación sistemática a gran escala de una zona de acceso público, como la videovigilancia del campus. Los supuestos (4), (6) y (9) reproducen los casos obligatorios del artículo~42 de la LOPDP, y los demás concretan, para el contexto universitario, la probabilidad de alto riesgo que el mismo artículo exige evaluar. La lista se implementa en la hoja \textit{EIPD previa} de la plantilla, que marca automáticamente los tratamientos con EIPD obligatoria. La regla RA-6 incorpora la consulta previa a la autoridad: la LOPDP no regula un procedimiento de consulta previa, pero faculta a la autoridad para requerir la EIPD y para pedir informes de riesgo \cite[arts.~9, lit.~c, y 42]{lopdp2021}, y la consulta antes de un tratamiento de alto riesgo residual es la práctica que recomiendan la guía de la AEPD y la norma ISO/IEC 29134 \cite{aepd2021riesgo,iso29134_2023}. Cada EIPD sigue la plantilla mínima del Anexo~\ref{anx:eipd}: descripción del tratamiento, necesidad y proporcionalidad, riesgos para los derechos de los titulares, medidas previstas y opinión del delegado.

Los plazos de las reglas RA-2 y RA-3 son parámetros de diseño propuestos por esta metodología; la institución que la adopte puede ajustarlos en la fase de preparación, siempre que lo documente.

\subsubsection{Estructura de la matriz de riesgos}

La matriz de riesgos se implementa como una hoja de cálculo con una fila por escenario y por dimensión (confidencialidad, integridad, disponibilidad o licitud). Sus campos se agrupan en seis bloques, que se resumen en la Tabla~\ref{tab:estructura_matriz}. Hojas adicionales contienen las escalas de probabilidad e impacto, la tabla de decisión, los criterios de aceptación, el catálogo de amenazas, el catálogo de controles referidos a ISO/IEC 27002:2022 \cite{iso27002_2022}, los resultados del panel Delphi, el cálculo del puntaje de Brier y las cinco hojas de la fase~0 y de la revisión previa: \textit{Normativa}, \textit{Perfil institucional}, \textit{Encargados y transferencias} ---con el proveedor, el servicio, los datos, el país, el contrato y la cláusula de protección de datos de cada encargado \cite[art.~41]{rlopdp2023}---, \textit{EIPD previa} y \textit{Retención}.

\begin{table}[htbp]
  \centering
  \caption{Bloques de campos de la matriz de riesgos.}
  \label{tab:estructura_matriz}
  \footnotesize
  \renewcommand{\arraystretch}{1.2}
  \begin{tabularx}{\textwidth}{|p{2.8cm}|L|}
    \hline
    \textbf{Bloque} & \textbf{Campos} \\
    \hline
    Identificación del tratamiento & ID, proceso, actividad de tratamiento, categorías de datos, categoría especial (sí/no), tipo y número estimado de titulares, grupo vulnerable, activo o sistema \\
    \hline
    Escenario de riesgo & Dimensión (C/I/D/licitud), comunidad de amenaza, vector, vulnerabilidad y su tipo (jurídica, organizacional, técnica, titulares), artículo de la LOPDP o norma de la SPDP, control existente \\
    \hline
    Riesgo inherente & Probabilidad anual y su justificación; impacto sobre el titular, impacto sobre el responsable y su justificación; agravantes aplicados; nivel inherente \\
    \hline
    Tratamiento & Opción (mitigar, transferir, evitar, aceptar), control propuesto con referencia a ISO/IEC 27002, responsable, plazo \\
    \hline
    Riesgo residual & Probabilidad residual, impacto residual, nivel residual, aceptación \\
    \hline
    Monitoreo & Indicador, frecuencia de revisión, estado, fecha de la última revisión \\
    \hline
  \end{tabularx}
  \raggedright
  \elaboradopor{Ramón Alejandro Manzaba Castro}
  \fuente{Elaboración propia}
\end{table}

El diseño, en conjunto, responde a cada requisito funcional con una actividad, un criterio o un campo de la matriz. Su rasgo distintivo frente a los modelos revisados en el Capítulo~\ref{chap:fundamentacion} es que el titular ocupa el centro de la valoración y que ninguna estimación queda sin justificación documentada.


\subsection{Catálogo de Escenarios Sectoriales}
\label{sec:escenarios_sectoriales}

La metodología incluye un catálogo de quince escenarios de riesgo propios del sector, que cada universidad usa como punto de partida y adapta a sus tratamientos. Cada escenario combina una amenaza documentada en la literatura sectorial \cite{ulven2021systematic,alvarez2024proteccion,li2023breach,mora2024lms} con el indicador del diagnóstico que lo activa. La Tabla~\ref{tab:escenarios} presenta los quince escenarios, cada uno referido a una sola dimensión.

\begin{table}[htbp]
  \centering
  \caption{Escenarios sectoriales de riesgo y sus indicadores activadores.}
  \label{tab:escenarios}
  \scriptsize
  \renewcommand{\arraystretch}{1.12}
  \begin{tabularx}{\textwidth}{|c|L|>{\raggedright\arraybackslash}p{2.05cm}|>{\raggedright\arraybackslash}p{1.8cm}|>{\raggedright\arraybackslash}p{2.1cm}|}
    \hline
    \textbf{ID} & \textbf{Escenario} & \textbf{Dimensión} & \textbf{Vulnerabilidad} & \textbf{Indicador activador} \\
    \hline
    R-01 & Exposición de certificados de discapacidad, datos de salud y fichas socioeconómicas por repositorios compartidos sin restricción & Confidencialidad & Organizacional & ID13; mapa de datos sensibles (D4) \\ \hline
    R-02 & Compromiso de cuentas del personal por \textit{phishing} y acceso a datos de estudiantes & Confidencialidad & Técnica & Registro de incidentes (\textit{phishing}) \\ \hline
    R-03 & Publicación de listados de admitidos, calificaciones o becarios con número de cédula, incluida la que se hace por transparencia activa (LOTAIP) sin aplicar minimización & Licitud & Organizacional & Revisión de publicaciones; perfil (pública) \\ \hline
    R-04 & Uso de datos para una finalidad no necesaria para el servicio educativo (difusión, publicidad) sin base válida o con un consentimiento que no es libre & Licitud & Jurídica & ID04, ID14 \\ \hline
    R-05 & Consultas indebidas al expediente por privilegios excesivos & Confidencialidad & Organizacional & Revisión de perfiles de acceso \\ \hline
    R-06 & Alteración de calificaciones mediante una cuenta privilegiada comprometida & Integridad & Técnica & Literatura sectorial \\ \hline
    R-07 & Tratamiento por proveedores de plataformas sin garantías contractuales & Confidencialidad & Jurídica & ID06a, ID17; inventario de encargados \\ \hline
    R-08 & Interrupción de los sistemas académicos por \textit{ransomware} & Disponibilidad & Técnica & Base obligatorio \\ \hline
    R-09 & Exfiltración de expedientes de becas en un ataque de doble extorsión & Confidencialidad & Técnica & ID13; RAT; antecedentes regionales \\ \hline
    R-10 & Conservación de datos sin una obligación legal ni una finalidad vigente & Licitud & Jurídica & ID05; tabla de retención \\ \hline
    R-11 & Transferencia internacional a servicios en la nube sin garantías evaluadas & Licitud & Jurídica & ID06b, ID17; inventario de encargados \\ \hline
    R-12 & Vulneración no notificada a la SPDP ni comunicada a los titulares en el término legal & Licitud & Organizacional & ID09; registro de incidentes \\ \hline
    R-13 & Derechos que no pueden ejercerse por falta de canal o por desconocimiento de los titulares & Licitud & Organizacional & ID07, ID08; registro de solicitudes \\ \hline
    R-14 & Supervisión automatizada de exámenes o analítica que elabora perfiles sin información previa ni garantía del derecho del art.~20 & Licitud & Jurídica & ID19, ID22; RAT \\ \hline
    R-15 & Extracción de datos personales y sensibles mediante un ataque informático y su publicación o venta & Confidencialidad & Técnica & Base obligatorio \\ \hline
  \end{tabularx}
  \raggedright
  \elaboradopor{Ramón Alejandro Manzaba Castro}
  \fuente{Elaboración propia con base en \cite{ulven2021systematic,alvarez2024proteccion,li2023breach,guerrero2021proctoring,mora2024lms}; \path{matriz/escenarios_sectoriales.json}}
\end{table}

Los escenarios cubren las cuatro dimensiones de análisis y tres de los cuatro tipos de vulnerabilidad de la guía de la SPDP \cite{spdp2026guiav2}: cinco son técnicos, cinco jurídicos y cinco organizacionales; la vulnerabilidad relativa a grupos vulnerables de titulares se incorpora mediante el agravante (a) en los escenarios con categorías especiales o datos sensibles (R-01, R-09, R-14 y R-15). Tres escenarios son \textbf{base obligatorios} y se analizan siempre, en toda universidad: R-06 (alteración de calificaciones), R-08 (interrupción por \textit{ransomware}) y R-15 (extracción de datos), porque dependen de controles técnicos ---autenticación, segmentación, respaldos, parches, registro--- que no pueden observarse desde fuera y, por tanto, ningún indicador externo puede descartarlos.

El \textbf{criterio de activación} de los demás escenarios lo aplica cada universidad con su propia fase~0: un escenario se activa cuando la lista de verificación califica su requisito activador como N o P, o cuando la evidencia interna ---el RAT, el inventario de encargados, el registro de incidentes o el registro de solicitudes de los titulares--- muestra el supuesto que lo origina. La institución puede ajustar ese criterio en la fase~P si documenta la justificación. Los requisitos de plataformas activan escenarios concretos: ID17 alimenta R-07, R-11 y el inventario de encargados; ID18, la revisión de la información al titular y de las bases de legitimación, e ID19, el escenario R-14 y el derecho del artículo~20, y si la supervisión usa biometría, la evaluación de impacto es obligatoria por la revisión previa y se aplica el agravante (a). La Tabla~\ref{tab:plan_tratamiento} vincula cada escenario con el requisito normativo que protege, los controles de referencia, el rol responsable, la regla que fija el plazo y el indicador de seguimiento.

\begin{table}[htbp]
  \centering
  \caption{Trazabilidad de los escenarios: requisito, control, responsable, plazo y seguimiento.}
  \label{tab:plan_tratamiento}
  \scriptsize
  \renewcommand{\arraystretch}{1.12}
  \begin{tabularx}{\textwidth}{|c|>{\raggedright\arraybackslash}p{2.3cm}|>{\raggedright\arraybackslash}p{2.4cm}|>{\raggedright\arraybackslash}p{2.0cm}|>{\raggedright\arraybackslash}p{1.3cm}|L|}
    \hline
    \textbf{ID} & \textbf{Requisito normativo} & \textbf{Controles (ISO/IEC 27002:2022)} & \textbf{Rol responsable tipo} & \textbf{Plazo} & \textbf{Indicador de seguimiento (meta)} \\
    \hline
    R-01 & LOPDP, arts.~25, 26 y 37 & 5.12; 5.15; 8.5; 8.12 & Bienestar estudiantil; TI & RA-2/RA-3 & Enlaces públicos en repositorios restringidos por mes (0) \\ \hline
    R-02 & LOPDP, art.~37 & 8.5; 6.3; 5.18 & TI & RA-2/RA-3 & Cuentas con autenticación multifactor (100\,\%) \\ \hline
    R-03 & LOPDP, arts.~7 y 10, lit.~e; LOTAIP, arts.~4, num.~5, y 9 (públicas) & 8.11 & Admisiones; secretaría & RA-2/RA-3 & Publicaciones con la cédula completa u otros identificadores innecesarios (0) \\ \hline
    R-04 & LOPDP, arts.~7, 8, 10, lit.~d, y 12 & 5.34; 5.31 & Comunicación; jurídico & RA-2/RA-3 & Finalidades con base de legitimación documentada (100\,\%) \\ \hline
    R-05 & LOPDP, art.~37 & 5.15; 5.18; 8.15 & Secretaría; TI & RA-2/RA-3 & Perfiles revisados al cierre del periodo (100\,\%) \\ \hline
    R-06 & LOPDP, arts.~10, lit.~h, y 37 & 8.2; 8.5; 8.15 & TI & RA-2/RA-3 & Cuentas privilegiadas con autenticación multifactor (100\,\%) \\ \hline
    R-07 & LOPDP, arts.~34 y 47, nums.~10 y 11; RLOPDP, art.~41 & 5.19; 5.20 & Jurídico; TI & RA-2/RA-3 & Encargados del inventario con contrato y cláusulas de protección de datos (100\,\%) \\ \hline
    R-08 & LOPDP, arts.~37 y 43 & 8.13; 8.8; 8.7; 5.24; 5.26 & TI & RA-2/RA-3 & Tiempo de restauración en la prueba de respaldo ($<24$~h) \\ \hline
    R-09 & LOPDP, arts.~25, 37, 43 y 46 & 8.24; 8.12; 8.16; 8.10; 8.5; 8.22 & TI; bienestar estudiantil & RA-2/RA-3 & Repositorios sensibles cifrados (100\,\%) \\ \hline
    R-10 & LOPDP, arts.~10, lit.~i, 12, num.~4, y 15 & 8.10; 5.9; 5.33 & Secretaría; archivo & RA-2/RA-3 & Expedientes sin obligación legal ni finalidad vigente que no se eliminaron ni anonimizaron (0) \\ \hline
    R-11 & LOPDP, arts.~12, num.~10, y 55 a 59 & 5.19; 5.20; 5.31 & TI; jurídico & RA-2/RA-3 & Transferencias internacionales con garantía documentada (100\,\%) \\ \hline
    R-12 & LOPDP, arts.~43, 46 y 68; Res.~SPDP-SPD-2026-0040-R & 5.24; 5.26 & Delegado o jurídico; TI & RA-2/RA-3 & Vulneraciones notificadas dentro del término (100\,\%) \\ \hline
    R-13 & LOPDP, arts.~12, nums.~14 a 16, y 13 a 17 & 5.34; 6.3 & Delegado o jurídico; secretaría & RA-2/RA-3 & Solicitudes atendidas en 15 días (100\,\%) \\ \hline
    R-14 & LOPDP, arts.~12, num.~17, 20, 26 y 42 & 5.34; 5.31 & Dirección académica; TI; jurídico & RA-2/RA-3 & Plataformas con EIPD, información previa y vía para impugnar decisiones automatizadas (100\,\%) \\ \hline
    R-15 & LOPDP, arts.~10, lit.~e, 25, 37, 43 y 46; Res.~SPDP-SPD-2026-0040-R & 8.5; 8.22; 8.2; 5.15; 8.15; 8.16; 8.12; 8.8; 8.10 & TI; delegado & RA-2/RA-3 & Sistemas con datos sensibles con MFA, segmentación y registro activo (100\,\%); expedientes vencidos suprimidos (100\,\%) \\ \hline
  \end{tabularx}
  \raggedright
  \elaboradopor{Ramón Alejandro Manzaba Castro}
  \fuente{Elaboración propia con base en \cite{lopdp2021,iso27002_2022,spdp2026vulneraciones}}
\end{table}

El plazo de cada escenario no se fija de antemano, sino que se deriva de su nivel de riesgo mediante las reglas RA-2 y RA-3; el plazo de 15 días del indicador de R-13 es el que la ley fija para atender las solicitudes de acceso y de rectificación \cite[arts.~13 y 14]{lopdp2021}. Los roles se expresan como funciones tipo, porque cada universidad que adopte la metodología los asignará según su estructura.

La construcción de escenarios a partir de requisitos tiene una ventaja que los catálogos genéricos no ofrecen: cada riesgo activado queda respaldado por una carencia verificable en la propia institución, y cada control propuesto responde a un artículo concreto de la ley.

\subsection{Fases 3 y 4: Respuesta a Vulneraciones y Monitoreo}
\label{sec:fase4_respuesta}

El procedimiento de respuesta a vulneraciones integra los plazos verificados en la Sección~\ref{sec:regimen_vulneraciones} y se estructura en seis pasos: (1) detección y registro del incidente, con la hora de detección, que es el inicio del cómputo de los plazos; (2) contención y evaluación inicial, en las primeras 24 horas, para determinar si hay datos personales afectados y si concurre alguna de las causales del artículo~24 del Reglamento General que hacen presumir un riesgo para los titulares \cite[art.~24]{rlopdp2023}; (3) clasificación de la vulneración por dimensión ---confidencialidad, integridad o disponibilidad--- y por su carácter temporal, permanente o irreversible, y valoración del riesgo para los titulares con las escalas de la metodología \cite{spdp2026vulneraciones}; (4) notificación a la SPDP a través del SISPDP, al organismo de regulación y control competente y, con fines informativos, al CSIRT, tan pronto sea posible y dentro del término de cinco días, con el contenido mínimo del artículo~26 del Reglamento General \cite[art.~43]{lopdp2021}, \cite[art.~26]{rlopdp2023}; si la universidad fuera calificada como operadora de infraestructura crítica digital, la notificación del incidente de ciberseguridad al ente rector se haría dentro de 72 horas \cite{lofc2026}; (5) comunicación a los titulares sin dilación y dentro del término de tres días, en lenguaje claro, salvo las excepciones del artículo~46 de la LOPDP \cite[art.~46]{lopdp2021}, y (6) análisis posterior, registro de lecciones aprendidas, actualización de la matriz y retroalimentación a la fase de identificación. La meta interna de la metodología es completar los pasos 1 a 4 en un máximo de 72 horas, de modo que el término legal funcione como límite y no como objetivo. La matriz se revisa, además, al cierre de cada periodo académico y ante cualquier cambio relevante de sistemas, proveedores o normativa.

\textbf{Caso de filtración publicada.} Cuando los datos extraídos ya se publicaron o se ofrecen a la venta (escenario R-15), el procedimiento se aplica con tres precisiones. Primera, la vulneración se presume de riesgo para los titulares, porque el responsable perdió el control de los datos y su tratamiento por terceros no está autorizado \cite[art.~24]{rlopdp2023}, de modo que la notificación a la SPDP se hace dentro del término de cinco días desde que se tuvo constancia, aunque la investigación técnica no haya terminado. Segunda, la comunicación a los titulares, dentro del término de tres días, debe informar la naturaleza del incidente, las categorías de datos expuestas, las fechas estimadas, las medidas adoptadas y recomendaciones concretas de protección ---cambiar contraseñas, activar la autenticación multifactor, vigilar intentos de suplantación o de fraude, y desconfiar de comunicaciones que usen los datos filtrados--- \cite[art.~46]{lopdp2021}, \cite{spdp2026vulneraciones}; si la comunicación individual exige un esfuerzo desproporcionado, procede una comunicación pública. Tercera, el análisis posterior revisa la minimización: los expedientes que no debían conservarse según la tabla de retención amplían el daño y deben suprimirse. Con esta fase, la metodología cierra su ciclo: lo que se aprende de cada incidente y de cada revisión vuelve a la identificación, y la gestión del riesgo deja de ser un ejercicio puntual para convertirse en un proceso continuo.


\subsection{Pruebas Técnicas de la Plantilla}
\label{sec:pruebas_tecnicas}

Esta sección y la siguiente distinguen dos tipos de evidencia que no deben confundirse. Las \textit{pruebas técnicas} comprueban que la plantilla calcula lo que el diseño dice que debe calcular; la \textit{validación metodológica} comprueba que la metodología es pertinente, clara, suficiente y coherente a juicio de expertos calibrados. La primera es condición necesaria de la segunda, pero no la sustituye.


Una metodología que delega sus cálculos en una hoja de cálculo solo es confiable si esas fórmulas son correctas. Por ello, las fórmulas de la plantilla se evaluaron con un motor de cálculo independiente (HyperFormula) en el momento de generar cada archivo: la plantilla y la matriz de escenarios sectoriales contienen 1\,150 fórmulas cada una y ninguna produjo errores en su última generación (25 de noviembre de 2026); las cuatro celdas que emplean la función \texttt{QUARTILE}, que ese motor no implementa, se verificaron con un cálculo equivalente en JavaScript. Además, se diseñaron seis casos de prueba con datos sintéticos para comprobar las reglas de decisión. La Tabla~\ref{tab:pruebas} muestra que todos produjeron el resultado esperado.

\begin{table}[htbp]
  \centering
  \caption{Casos de prueba de las reglas de decisión de la plantilla.}
  \label{tab:pruebas}
  \scriptsize
  \renewcommand{\arraystretch}{1.2}
  \begin{tabularx}{\textwidth}{|c|L|>{\raggedright\arraybackslash}p{3.2cm}|>{\raggedright\arraybackslash}p{3.2cm}|c|}
    \hline
    \textbf{Caso} & \textbf{Regla verificada} & \textbf{Entrada} & \textbf{Resultado obtenido} & \textbf{Estado} \\
    \hline
    T-01 & Impacto muy alto con probabilidad muy baja no desciende de nivel alto & P = Muy baja; I = Muy alto & Nivel inherente: Alto & Correcto \\
    \hline
    T-02 & Probabilidad muy alta con impacto muy bajo no se trata como riesgo alto & P = Muy alta; I = Muy bajo & Nivel inherente: Medio & Correcto \\
    \hline
    T-03 & Los agravantes no superan el nivel muy alto & I = Alto; 3 agravantes & Impacto ajustado: Muy alto & Correcto \\
    \hline
    T-04 & Una fila incompleta no genera nivel ni errores & Sin probabilidad & Celdas de resultado vacías & Correcto \\
    \hline
    T-05 & Riesgo residual alto sin categoría especial no es aceptable & PR = Media; IR = Alto & Nivel residual: Alto; ``No aceptable: tratamiento prioritario'' & Correcto \\
    \hline
    T-06 & Impacto residual muy alto exige aprobación reforzada & PR = Muy baja; IR = Muy alto & Nivel residual: Alto; requiere DPD/jurídico y máxima autoridad & Correcto \\
    \hline
  \end{tabularx}
  \raggedright
  \elaboradopor{Ramón Alejandro Manzaba Castro}
  \fuente{Ejecución de \texttt{scripts/build\_matriz.js} con los casos de prueba}
\end{table}

Las pruebas T-01 y T-06 tienen un significado que va más allá de lo técnico: demuestran que la herramienta impide, por construcción, que un daño muy grave para el titular quede oculto detrás de una probabilidad baja, que es la regla de la guía de la SPDP que los modelos basados en la multiplicación de etiquetas no garantizan.


Los seis casos de prueba se ejecutaron sobre la versión de la plantilla de octubre de 2026; la generalización de noviembre de 2026, que eliminó las referencias a un proceso institucional específico, modificó únicamente los rótulos, las listas de procesos y los catálogos de amenazas y controles. La versión del 25 de noviembre de 2026 añadió cinco hojas ---\textit{Normativa}, \textit{Perfil institucional}, \textit{Encargados y transferencias}, \textit{EIPD previa} y \textit{Retención}--- con 93 fórmulas nuevas, sin alterar las escalas, la tabla de decisión ni las fórmulas existentes; las 1\,150 fórmulas se evaluaron sin errores.

\subsection{Justificación Técnica}

La selección de los componentes de la metodología responde a seis decisiones. \textbf{Marco de proceso}: se adoptan las etapas de la guía de la SPDP, alineadas con ISO/IEC 27005:2022, porque la guía es la referencia normativa vigente en Ecuador y su capítulo~I es de cumplimiento obligatorio \cite{spdp2025guia,spdp2026guiav2}; se descartó adoptar un marco extranjero completo, como el de la AEPD o el NIST Privacy Framework, porque ninguno responde a la LOPDP \cite{aepd2021riesgo,nist2020privacy}. \textbf{Enfoque centrado en el titular}: a diferencia de los modelos aplicados en el sector público y educativo ecuatoriano, centrados en activos \cite{patino2018ict,lopezvasco2025iso}, la escala de impacto principal se refiere al daño al titular, como respalda la literatura sobre categorías de daño a la privacidad \cite{wairimu2025harm}. \textbf{Escenarios activados por evidencia}: los escenarios se activan por carencias verificables en la propia institución y no por la intuición del analista. \textbf{Buenas prácticas comparadas}: los elementos que la LOPDP no detalla ---la lista de supuestos de evaluación de impacto, la consulta previa, el registro de actividades o la escala de madurez--- se toman de marcos que ya los aplican, sin exigir una reforma legal. \textbf{Estimación ordinal calibrada}: los modelos cuantitativos como FAIR requieren datos de frecuencia y pérdidas que una universidad rara vez posee \cite{wang2020fair}; por ello se opta por escalas ordinales ancladas en frecuencias anuales, estimadas mediante un panel Delphi y validadas con el puntaje de Brier, opción que la guía de la SPDP admite cuando se documenta la justificación de cada valor \cite{spdp2026guiav2}. \textbf{Herramienta de soporte}: la matriz se implementa en una hoja de cálculo compatible con Microsoft Excel y LibreOffice Calc, lo que reduce el costo de adopción y cumple los requisitos RNF02, RNF05 y RNF07.

Podría objetarse que una metodología construida con fuentes documentales y validada por expertos, y no aplicada a una institución real, carece de prueba de utilidad. La objeción se responde en dos niveles. En el plano del diseño, la metodología no se apoya en supuestos sobre una institución, sino en las obligaciones de la LOPDP y en soluciones que otros marcos ya aplican, con escenarios trazables a la ley; por eso es aplicable a cualquier universidad que reproduzca la fase~0 con sus propios datos. En el plano de la evidencia, el trabajo no afirma que la metodología haya reducido el riesgo de ninguna institución, sino que está completamente especificada, que su herramienta es correcta y que expertos calibrados la juzgan válida según criterios explícitos. En conclusión, la propuesta convierte el análisis comparado en una herramienta de gestión que cualquier universidad puede adoptar, y declara con precisión qué parte de su validez está demostrada y qué parte depende de la evaluación del panel.
